% Problem Set 7 -- 6 problems = Munkres, Topology (2nd ed.), Section 22 Exercises 1-6.
% Do not edit except to fix statement typos.
\problemset{7}{Munkres \S22}

\problem{7}{1}{Check the details of Example 3.\footnote{Example 3 of \S22 (The Quotient
Topology) in J.\,R.~Munkres, \emph{Topology}, 2nd ed.}}

\problem{7}{2}{\leavevmode
\begin{enumerate}[label=(\alph*)]
  \item Let $p\colon X\to Y$ be a continuous map. Show that if there is a continuous map
        $f\colon Y\to X$ such that $p\circ f$ equals the identity map of $Y$, then $p$ is a
        quotient map.
  \item If $A\subset X$, a \emph{retraction} of $X$ onto $A$ is a continuous map
        $r\colon X\to A$ such that $r(a)=a$ for each $a\in A$. Show that a retraction is a
        quotient map.
\end{enumerate}}

\problem{7}{3}{Let $\pi_1\colon\R\times\R\to\R$ be projection on the first coordinate. Let
$A$ be the subspace of $\R\times\R$ consisting of all points $x\times y$ for which either
$x\ge 0$ or $y=0$ (or both); let $q\colon A\to\R$ be obtained by restricting $\pi_1$. Show
that $q$ is a quotient map that is neither open nor closed.}

\problem{7}{4}{\leavevmode
\begin{enumerate}[label=(\alph*)]
  \item Define an equivalence relation on the plane $X=\R^2$ as follows:
        \[ x_0\times y_0\sim x_1\times y_1\quad\text{if}\quad x_0+y_0^2=x_1+y_1^2. \]
        Let $X^*$ be the corresponding quotient space. It is homeomorphic to a familiar
        space; what is it? [\emph{Hint:} Set $g(x\times y)=x+y^2$.]
  \item Repeat (a) for the equivalence relation
        \[ x_0\times y_0\sim x_1\times y_1\quad\text{if}\quad x_0^2+y_0^2=x_1^2+y_1^2. \]
\end{enumerate}}

\problem{7}{5}{Let $p\colon X\to Y$ be an open map. Show that if $A$ is open in $X$, then the
map $q\colon A\to p(A)$ obtained by restricting $p$ is an open map.}

\problem{7}{6}{Recall that $\R_K$ denotes the real line in the $K$-topology (see \S13). Let
$Y$ be the quotient space obtained from $\R_K$ by collapsing the set $K$ to a point; let
$p\colon\R_K\to Y$ be the quotient map.
\begin{enumerate}[label=(\alph*)]
  \item Show that $Y$ satisfies the $T_1$ axiom, but is not Hausdorff.
  \item Show that $p\times p\colon\R_K\times\R_K\to Y\times Y$ is not a quotient map.
        [\emph{Hint:} The diagonal is not closed in $Y\times Y$, but its inverse image is
        closed in $\R_K\times\R_K$.]
\end{enumerate}}
